
%(BEGIN_QUESTION)
% Copyright 2006, Tony R. Kuphaldt, released under the Creative Commons Attribution License (v 1.0)
% This means you may do almost anything with this work of mine, so long as you give me proper credit

Les og lag en disposisjon av delkapitlet ``Electronic Positioners'' i avsnittet ``Valve Positioners'' i kapitlet ``Control Valves'' i læreboken {\it Lessons In Industrial Instrumentation}. Noter sidetallene der viktige illustrasjoner, fotografier, ligninger, tabeller og andre relevante detaljer finnes. Forbered deg på å diskutere begrepene og eksemplene fra denne lesingen grundig med lærer og medstudenter.

\underbar{file i01365}
%(END_QUESTION)





%(BEGIN_ANSWER)

 
%(END_ANSWER)





%(BEGIN_NOTES)

Denne typen posisjonerer måler ventilspindelens posisjon elektronisk, og påfører deretter lufttrykk for å oppnå målposisjonen. Selvdiagnostikk er en ekstra fordel.

\vskip 10pt

Ved svikt i posisjonssensoren går noen smarte posisjonerere over i volume-booster-modus: de videresender ganske enkelt omtrentlig trykkmengde som trengs for å posisjonere ventilen for et gitt inngangssignal, basert på tidligere trykk-/posisjonsdata.

\vskip 10pt

Diagnosegrafen ``valve signature'': en kurve for aktuatortrykk (kraft) mot spindelbevegelse, både ved åpning og lukking. Avstanden mellom åpnings- og lukkekurvene representerer friksjonen i ventilmekanismen. Jo større friksjon, desto større avstand mellom kurvene. Nedre venstre del av grafen er seteprofilen: en registrering av kraft mot bevegelse når pluggen kommer i kontakt med setet. Ventilsignaturer kan registreres og arkiveres for senere sammenligning når ventilen slites. Elektriske ventilaktuatorer kan registrere tilsvarende signaturer ved å plotte motorstrøm mot spindelposisjon.


\vskip 10pt

Kanskje den største fordelen med digitale ventilposisjonerere er {\it diagnosekapasiteten} de tilfører reguleringsventiler. Ved å måle hvor mye trykk som trengs for å bevege en ventilaktuator ved en gitt endring i inngangssignal, kan en digital posisjoner indirekte måle pakkboksfriksjon, friksjon mellom plugg og sete, og andre parametere som indikerer ventilens helse.






\vskip 20pt \vbox{\hrule \hbox{\strut \vrule{} {\bf Forslag til sokratisk diskusjon} \vrule} \hrule}

\begin{itemize}
\item{} Identifiser viktige funksjoner i digitale, elektroniske ventilposisjonerere som mekaniske posisjonerere mangler.
\item{} Studer blokkdiagrammet av en elektronisk posisjoner i læreboken, og forklar hvordan de to reguleringssløyfene fungerer for å holde ventilspindelen i konstant posisjon (ved konstant reguleringssignal) til tross for ytre påvirkninger på spindelen (f.eks. pakkboksfriksjon, krefter på pluggen fra prosessstrømning, endringer i luftforsyningstrykk).
\item{} Studer blokkdiagrammet av en elektronisk posisjoner i læreboken, og kommenter ``+''- og ``$-$''-symbolene ved hver summeringsblokk ($\Sigma$). Forklar hvordan disse symbolene tilsvarer ikke-inverterende og inverterende innganger på en operasjonsforsterker.
\item{} Forklar hvilke typer nyttige data som vises i en ``valve signature''-graf.
\item{} Forklar hvorfor endene av ``valve signature''-grafen i boken går brått vertikalt.
\item{} Forklar hvordan en ``valve signature''-graf kan avdekke slitasje i ventilsetet.
\end{itemize}

%INDEX% Final Control Elements, valve: positioner (digital)

%(END_NOTES)
